\section{Experimental Setup}

\subsection{Software environment}

The repository uses the Qiskit software stack \cite{qiskit} and pins Qiskit 2.5.1, Qiskit Aer 0.17.2, Qiskit IBM Runtime 0.48.0, NumPy 2.x, pandas 2.x, SciPy 1.14 or newer, Matplotlib 3.8 or newer, and NetworkX 3.2 or newer. The notebook is designed for sequential execution in Google Colab or an equivalent Python environment. COBYLA is provided through SciPy \cite{virtanen2020}.

The random seed for data generation, statevector sampling, and transpilation is 2026. Standard-\QAOA{} candidate tables use 4096 ideal shots. The hardware comparison uses 1024 shots per circuit.

\subsection{Ideal depth protocol}

Standard \QAOA{} is optimized for
\begin{equation}
p\in\{1,2,3\}
\end{equation}
on all three instances. The original $3\times3$, $p=1$ experiment is retained. No additional parameter grid, multiple optimization seeds, or depth values are introduced. \GMQAOA{} is optimized at $p=1$ for each fixed instance.

\subsection{IBM Quantum execution}

The physical experiment targets the $4\times2$ exact-capacity instance and uses $p=1$ for both algorithms. The archived job metadata are summarized in Table~\ref{tab:job_metadata}.

\begin{table}[H]
\centering
\caption{IBM Quantum job metadata.}
\label{tab:job_metadata}
\begin{tabular}{ll}
\toprule
Field & Recorded value\\
\midrule
Backend & \texttt{ibm\_marrakesh}\\
Backend qubits & 156\\
Job identifier & \texttt{d9vl8qfo3ppc73akfak0}\\
Created (UTC) & 2026-08-14 17:46:17.684431\\
Shots per circuit & 1024\\
Circuits in job & 2\\
Twirling realizations & 16 per circuit\\
Shots per realization & 64\\
Dynamical decoupling & Enabled, XpXm\\
Gate twirling & Enabled\\
Estimated running time & \SI{5.075875}{s}\\
Recorded execution span & \SI{4.400754}{s}\\
Client wall time & \SI{13.695772}{s}\\
\bottomrule
\end{tabular}
\end{table}
The execution span is derived from Runtime timestamps. Client wall time includes submission, service communication, and result retrieval and is not an algorithm-specific QPU duration.

\subsection{Transpilation}

Both measured circuits use the same backend-specific preset pass manager:
\begin{lstlisting}[style=pythonstyle,caption={Backend-specific compilation used in the notebook.}]
pass_manager = generate_preset_pass_manager(
    backend=target_backend,
    optimization_level=3,
    seed_transpiler=2026,
)
standard_isa = pass_manager.run(standard_measured)
gm_isa = pass_manager.run(gm_measured)
\end{lstlisting}
This process selects a physical layout, decomposes non-native gates, and routes interactions through the backend coupling graph. The final native two-qubit operation in the recorded circuits is CZ; the exported CX and ECR counts are zero.

Figure~\ref{fig:topology_overhead} summarizes the logical-to-physical mismatch visible in the executed notebook.
\begin{figure}[H]
\centering
\includegraphics[width=0.98\textwidth]{transpilation_topology_overhead.png}
\caption{Topology and transpilation-overhead summary generated by the executed notebook. Dense logical assignment couplings must be routed on a sparse device graph.}
\label{fig:topology_overhead}
\end{figure}

\subsection{Distribution and uncertainty metrics}

For ideal distribution $p$ and measured distribution $q$,
\begin{align}
\TVD(p,q)&=\frac12\sum_x|p(x)-q(x)|,\label{eq:tvd}\\
\Fcl(p,q)&=\left(\sum_x\sqrt{p(x)q(x)}\right)^2,\label{eq:fidelity}\\
H(p,q)&=\sqrt{1-\sum_x\sqrt{p(x)q(x)}}.
\label{eq:hellinger}
\end{align}
$\Fcl$ is a classical fidelity between measurement distributions, not full quantum-state fidelity on the device. Feasible probability, optimal probability, and leakage are
\begin{equation}
\PF=\sum_{x\in\mathcal F}q(x),
\qquad
\Popt=\sum_{x:f(x)=f^\star}q(x),
\qquad
L_{\mathcal F}=1-\PF.
\end{equation}
Wilson 95\% intervals are used for binomial feasible and optimal counts \cite{wilson1927}.

\subsection{Reconstruction of the complete analysis distribution}

The archived Runtime job contains the complete 1024-shot BitArray for each circuit, from which the full counts are decoded. The \GMQAOA{} ideal distribution is exactly reproduced from its exported $\beta$ and $\gamma$. The repository exports standard-\QAOA{} ideal probabilities to six decimal places but does not retain the full ideal distribution as a separate file. For global standard-\QAOA{} distances, the analytical $p=1$ state is reconstructed by matching those exported probabilities and the exported expected energy, feasible probability, and optimal probability. The recovered values are $\beta\approx1.8352733$ and $\gamma\approx1.2267104$. Standard global distances are therefore reported to four significant digits and identified as reconstructed within the repository's export precision.
